\subsection{Local Non-Affine T1 Observable}

Let $\mathbf{x}_i(t)$ and $\mathbf{v}_i(t)$ denote the position and velocity of
individual $i$ at time $t$. For each focal individual, a local neighborhood was
defined geometrically. The local affine component of motion was estimated from
the neighborhood and subtracted from the focal-centered motion. The remaining
component is the local non-affine residual.

This paper focuses on the tangential component of that residual around
transition events in a compact-density coordinate. We refer to this frozen
measurement target as T1. T1 is therefore not an assumed force, behavioral
rule, or complete mechanism. It is a transition-linked local tangential
non-affine observable:
\begin{equation}
T1(t)
=
\mathrm{Tangential}
\left[
\mathbf{v}_{\mathrm{local}}(t)
-
\mathbf{v}_{\mathrm{affine}}(t)
\right],
\end{equation}
evaluated relative to event-conditioned windows and matched controls. The
important point is conceptual rather than notational: T1 is measured only after
local affine deformation has been removed.

